The \code{BPatch\_basicBlock} class represents a basic block in the application being
instrumented. Objects of this class representing the blocks within a function
can be obtained using the \code{BPatch\_flowGraph} object for the function.
\code{BPatch\_basicBlock} includes methods for navigating through the control flow
graph of the containing function.

\begin{apient}
  void getSources(std::vector<BPatch_basicBlock*>&)
\end{apient}

\apidesc{
  Fills the given vector with the list of predecessors for this basic block (i.e,
  basic blocks that have an outgoing edge in the control flow graph leading to
  this block).
}

\begin{apient}
  void getTargets(std::vector<BPatch_basicBlock*>&)
\end{apient}

\apidesc{
  Fills the given vector with the list of successors for this basic block (i.e,
  basic blocks that are the destinations of outgoing edges from this block in the
  control flow graph).
}

\begin{apient}
  void getOutgoingEdges(std::vector<BPatch_edge *> &out)
\end{apient}

\apidesc{
  Fill out with all of the control flow edges that leave this basic block.
}

\begin{apient}
  void getIncomingEdges(std::vector<BPatch_edge *> &inc)
\end{apient}

\apidesc{
  Fills inc with all of the control flow edges that point to this basic block.
}

\begin{apient}
  bool getInstructions(std::vector<Dyninst::InstructionAPI::Instruction>&)
  
  bool getInstructions(std::vector<std::pair<Dyninst::InstructionAPI::Instruction, Address> >&)
\end{apient}

\apidesc{
  Fills the given vector with InstructionAPI Instruction objects representing the
  instructions in this basic block, and call also returns the address each instruction
  starts at.
}

\begin{apient}
  bool dominates(BPatch_basicBlock*)
\end{apient}

\apidesc{
  This function returns true if the argument is pre-dominated in the control flow
  graph by this block, and false if it is not.
}

\begin{apient}
  BPatch_basicBlock* getImmediateDominator()
\end{apient}

\apidesc{
  Return the basic block that immediately pre-dominates this block in the control flow graph.
}

\begin{apient}
  void getImmediateDominates(std::vector<BPatch_basicBlock*>&)
\end{apient}

\apidesc{
  Fill the given vector with a list of pointers to the basic blocks that are
  immediately dominated by this basic block in the control flow graph.
}
\begin{apient}
  void getAllDominates(std::set<BPatch_basicBlock*>&)
  void getAllDominates(BPatch_Set<BPatch_basicBlock*>&)
\end{apient}

\apidesc{
  Fill the given set with pointers to all basic blocks that are dominated by this
  basic block in the control flow graph.
}

\begin{apient}
  bool getSourceBlocks(std::vector<BPatch_sourceBlock*>&)
\end{apient}

\apidesc{
  Fill the given vector with pointers to the source blocks contributing to this
  basic block\'s instruction sequence.
}
\begin{apient}
  int getBlockNumber()
\end{apient}

\apidesc{
  Return the ID number of this basic block. The ID numbers are consecutive from 0 to
  n-1, where n is the number of basic blocks in the flow graph to which this basic
  block belongs.
}
\begin{apient}
  std::vector<BPatch_point *> findPoint(const std::set<BPatch_opCode> &ops)
  std::vector<BPatch_point *> findPoint(const BPatch_Set<BPatch_opCode> &ops)
\end{apient}

\apidesc{
  Find all points in the basic block that match the given operation.
}
\begin{apient}
  BPatch_point *findEntryPoint(
  BPatch_point *findExitPoint()
\end{apient}

\apidesc{
  Find the entry or exit point of the block.
}

\begin{apient}
  unsigned long getStartAddress()
\end{apient}

\apidesc{
  This function returns the starting address of the basic block. The address returned
  is an absolute address.
}
\begin{apient}
  unsigned long getEndAddress()
\end{apient}

\apidesc{
  This function returns the end address of the basic block. The address returned is
  an absolute address.
}
\begin{apient}
  unsigned long getLastInsnAddress()
\end{apient}

\apidesc{
  Return the address of the last instruction in a basic block.
}
\begin{apient}
  bool isEntryBlock()
\end{apient}

\apidesc{
  This function returns true if this basic block is an entry block into a function.
}
\begin{apient}
  bool isExitBlock()
\end{apient}

\apidesc{
  This function returns true if this basic block is an exit block of a function.
}
\begin{apient}
  unsigned size()
\end{apient}

\apidesc{
  Return the size of a basic block. The size is defined as the difference between the
  end address and the start address of the basic block.
}
